\ticket{8}{Онтологии в инженерии знаний}

\begin{examplan}
    \item Дать понятие онтологии.
    \item Описать состав онтологии и типы отношений концептов.
    \item Перечислить классификации онтологий и привести примеры.
    \item Назвать принципы и основные шаги построения онтологии.
\end{examplan}

\subsection*{Короткая версия для письменного ответа}

\textbf{Онтология} в инженерии знаний --- формальное описание предметной области, включающее концепты, отношения, свойства, аксиомы и экземпляры. Онтология задает общий словарь и правила интерпретации понятий, чтобы знания могли использоваться людьми и программными системами.

\subsection*{Состав онтологии}

\begin{itemize}
    \item \textbf{Концепты} или классы: основные понятия предметной области.
    \item \textbf{Экземпляры}: конкретные объекты, принадлежащие концептам.
    \item \textbf{Свойства данных}: атрибуты, связывающие объект с литеральным значением.
    \item \textbf{Объектные свойства}: отношения между экземплярами классов.
    \item \textbf{Аксиомы}: ограничения и утверждения, всегда истинные в рамках онтологии.
    \item \textbf{Правила и ограничения}: домены, диапазоны, кардинальность, транзитивность, симметричность.
\end{itemize}

\subsection*{Типы отношений концептов}

\begin{itemize}
    \item \textbf{IS-A} или \textbf{subClassOf}: отношение вид-род, например $Mammal\sqsubseteq Animal$.
    \item \textbf{instanceOf}: принадлежность индивида классу.
    \item \textbf{partOf/hasPart}: часть-целое.
    \item \textbf{Ассоциативные отношения}: причинность, лечение, использование, расположение и другие предметные связи.
    \item \textbf{Атрибутивные отношения}: связь объекта с характеристикой, например возрастом или датой.
\end{itemize}

\subsection*{Классификация онтологий}

Онтологии классифицируют по разным основаниям:
\begin{itemize}
    \item по степени формальности: неформальные, полуформальные, формальные;
    \item по глубине: легковесные без развитой системы аксиом и тяжеловесные с аксиомами;
    \item по охвату: верхнего уровня, предметные, прикладные;
    \item по методологии: лингвистические, строящиеся от слов к понятиям, и понятийные, строящиеся от концептов;
    \item по типу отношений: таксономии, партономии, атрибутивные, причинно-следственные, неоднородные.
\end{itemize}

\subsection*{Примеры}

\begin{itemize}
    \item \textbf{WordNet} --- лексический ресурс и легковесная лингвистическая онтология.
    \item \textbf{Cyc} --- формальная онтология общих знаний и здравого смысла.
    \item \textbf{SUMO} --- онтология верхнего уровня.
    \item Доменные медицинские или биологические онтологии описывают специализированные предметные области.
    \item Рубрикаторы и классификаторы могут рассматриваться как прикладные легковесные онтологические ресурсы.
\end{itemize}

\subsection*{Принципы построения}

\begin{itemize}
    \item \textbf{Ясность}: определения должны быть понятны и недвусмысленны.
    \item \textbf{Непротиворечивость}: онтология не должна давать противоречивые выводы.
    \item \textbf{Расширяемость}: новые понятия должны добавляться без разрушения существующей структуры.
    \item \textbf{Минимальность}: не стоит вводить лишние сущности без необходимости.
    \item \textbf{Переиспользование}: полезно опираться на уже существующие онтологии и стандарты.
\end{itemize}

\subsection*{Основные шаги построения}

\begin{enumerate}
    \item Определить предметную область, цель и пользователей онтологии.
    \item Сформулировать компетентностные вопросы, на которые онтология должна помогать отвечать.
    \item Найти и оценить существующие онтологии для переиспользования.
    \item Выделить ключевые термины и концепты.
    \item Построить иерархию классов: сверху вниз, снизу вверх или комбинированно.
    \item Определить свойства, отношения, домены, диапазоны и ограничения.
    \item Создать экземпляры, если они входят в задачу.
    \item Закодировать онтологию, например в RDF Schema или OWL.
    \item Проверить полноту, непротиворечивость и соответствие целям.
    \item Документировать и сопровождать онтологию.
\end{enumerate}

\begin{important}
    \item Онтология задает формальный словарь предметной области.
    \item В состав онтологии входят концепты, отношения, свойства, аксиомы и экземпляры.
    \item Важнейшие отношения: IS-A, instanceOf, partOf и предметные ассоциации.
    \item Верхнеуровневые онтологии описывают общие категории, доменные --- конкретную область, прикладные --- задачу.
    \item Построение онтологии начинается с цели и компетентностных вопросов.
\end{important}

